\subsection{Resistor:}
A resistor is a two-port passive device that opposes the flow of current,
we call this property resistance. The unit of resistance is Ohm (\(\Omega\)).
Everything has resistance. A good way to imagine resistance is to
think about it's value as a constant of proportionality between voltage
and current. The relationship is given by Ohm's law:
\[
  V = I_ \cdot R
\]
The inverse of resistance is called conductance (\(G\)) and is measured
in Siemens (S). The relationship between resistance and conductance is:
\[
  G = \frac{1}{R}
\]

You'll notice when we apply KCL to circuits within this book, we'll mostly
use conductance \(G\) to express the resistance, we do this as it
simplifies the analysis a lot. The conductance version of Ohm's law is:
\[
  I = V \cdot G
\]
An ideal resistor is a linear (i.e. the r1elationship between voltage and
current is a straight line) and time-invariant (i.e. the relationship
between voltage and current does not change with time) of the three passive
components. It is also the only device that in its ideal form does not
store energy. This is why at low frequencies, if a resistor can be used
in place of a capacitor or inductor, use it.

\subsection{Capacitor:}
A capacitor is a two-port passive device that stores energy in the form of
an electric field. The property of a capacitor is called capacitance and is
measured in Farads (F). The core relationship between charge, capacitance
and voltage is given by:
\begin{align*}
  Q = C \cdot V \\
  \frac{dQ}{dt} &= C \cdot \frac{dV}{dt} \hspace{1cm} \text{(assuming \(C\) is constant)} \\
  I &= C \cdot \frac{dV}{dt} \hspace{1cm} [I = \frac{dQ}{dt}]
\end{align*}
This is the fundamental relationship between current and voltage in a
capacitor or if you may call this the capacitor's version of Ohm's law.
The inverse of capacitance is called elastance (\(S\)) and is measured in Farads
inverse (F\(^{-1}\)).

The impedance of a capacitor is given by:
\[
  Z_C = \frac{1}{j\omega C} = \frac{1}{j 2 \pi f C} = \frac{1}{s C}
\]
The inverse of impedance is called admittance (\(Y\)) and is measured in
Siemens (S). The relationship between impedance and admittance is:
\[
  Y = \frac{1}{Z} = \frac{1}{\frac{1}{sC}} = sC
\]
You'll notice when we apply KCL to circuits within this book, we'll mostly
use admittance \(Y\) to express the impedance, we do this as it simplifies
the analysis a lot. The admittance version of Ohm's law is:
\[
  I = V \cdot Y
\]
The impedance of a capacitor is inversely proportional to frequency. This
means that at low frequencies, the impedance of a capacitor is very high and
at frequencies, the impedance of a capacitor is very low. This is why
capacitor effects are more pronounced at high frequencies as their
impedance drops enough to load the circuit. This is also why capacitors are
used for decoupling (they act as a short circuit at high frequencies).

An ideal capacitor is a linear (i.e. the relationship between voltage and
current is a straight line) and time-invariant (i.e. the relationship
between voltage and current does not change with time) of the three passive
components. It is also the only device that in its ideal form does not
dissipate energy.

A capacitor can be formed anywhere two conductors are separated by an
insulator. The charges are stored on the surface of the  conductors and the
energy is stored in the electric field between the  conductors. The energy
stored in a capacitor is given by:
\begin{align*}
  I &= C \cdot \frac{dV}{dt} \\
  P &= V \cdot I \\
  P &= V \cdot C \cdot \frac{dV}{dt} \\
  P &= C \cdot V \cdot \frac{dV}{dt} \\
  P &= C \cdot \frac{1}{2} \cdot \frac{d(V^2)}{dt} \hspace{1cm} [\frac{d(V^2)}{dt} = 2V \cdot \frac{dV}{dt}] \\
  \int P dt &= \int C \cdot \frac{1}{2} \cdot \frac{d(V^2)}{dt} dt \\
  E &= \frac{1}{2} C V^2
\end{align*}

\subsection{Inductor:}
An inductor is a two-port passive device that stores energy in the form of
a magnetic field. The property of an inductor is called inductance and is
measured in Henrys (H). The core relationship between flux, inductance
and current is given by:
\begin{align*}
  \Phi &= L \cdot I \\
  \frac{d\Phi}{dt} &= L \cdot \frac{dI}{dt} \hspace{1cm} \text{(assuming \(L\) is constant)} \\
  V &= L \cdot \frac{dI}{dt} \hspace{1cm} [V = \frac{d\Phi}{dt}]
\end{align*}
This is the fundamental relationship between voltage and current in an
inductor or if you may call this the inductor's version of Ohm's law.
The inverse of inductance is called reluctance (\(R\)) and is measured in
Henrys inverse (H\(^{-1}\)).
\[
  R = \frac{1}{L}
\]
The impedance of an inductor is given by:
\[
  Z_L = j\omega L = j 2 \pi f L = s L
\]
The inverse of impedance is called admittance (\(Y\)) and is measured in
Siemens (S). The relationship between impedance and admittance is:
\[
  Y = \frac{1}{Z} = \frac{1}{sL}
\]
You'll notice when we apply KCL to circuits within this book, we'll mostly
use admittance \(Y\) to express the impedance, we do this as it simplifies
the analysis a lot. The admittance version of Ohm's law is:
\[
  I = V \cdot Y
\]
The impedance of an inductor is directly proportional to frequency. At low
frequencies, the impedance of an inductor is very low and at high frequencies,
the impedance of an inductor is very high. This is why inductor effects are
more pronounced at low frequencies as their impedance rises enough to load
the circuit. This is also why inductors are used for filtering (they act as
a short circuit at low frequencies).

An ideal inductor is a linear (i.e. the relationship between voltage and
current is a straight line) and time-invariant (i.e. the relationship
between voltage and current does not change with time) of the three passive
components. It is also the only device that in its ideal form does not
dissipate energy. The energy stored in an inductor is given by:
\begin{align*}
  V &= L \cdot \frac{dI}{dt} \\
  P &= V \cdot I \\
  P &= L \cdot \frac{dI}{dt} \cdot I \\
  P &= L \cdot I \cdot \frac{dI}{dt} \\
  P &= L \cdot \frac{1}{2} \cdot \frac{d(I^2)}{dt} \hspace{1cm} [\frac{d(I^2)}{dt} = 2I \cdot \frac{dI}{dt}] \\
  \int P dt &= \int L \cdot \frac{1}{2} \cdot \frac{d(I^2)}{dt} dt \\
  E &= \frac{1}{2} L I^2
\end{align*}

